% Generated draft: requires mathematical and source review.
\documentclass[11pt]{article}
\usepackage[margin=25mm]{geometry}
\usepackage{fontspec}
\setmainfont{Latin Modern Roman}
\usepackage{amsmath,amssymb,mathtools,unicode-math}
\setmathfont{Latin Modern Math}
\usepackage{enumitem,needspace,xurl,hyperref}
\setlength{\parindent}{0pt}
\setlength{\parskip}{5pt}
\setlength{\emergencystretch}{3em}
\newcommand{\N}{\mathbb{N}}
\newcommand{\Z}{\mathbb{Z}}
\newcommand{\Q}{\mathbb{Q}}
\newcommand{\R}{\mathbb{R}}
\newcommand{\C}{\mathbb{C}}
\newcommand{\F}{\mathbb{F}}
\newcommand{\A}{\mathbb{A}}
\newcommand{\PP}{\mathbb{P}}
\newcommand{\E}{\mathbb{E}}
\newcommand{\eps}{\varepsilon}
\newcommand{\dd}{\,\mathrm{d}}
\begin{document}
\section*{UnsolvedMath ID: 30006429}
\textbf{Source problem code:} catalogue entry
\subsection{Definitions and mathematical statement}
The supplied question concerns subalgebra zeta functions and their proposed relationship with flag Hilbert--Poincar\'e series. For the combinatorial context, let \(K\) be a number field and let \(\mathcal A\) be a finite set of hyperplanes in affine space \(\mathbb A^n_K\). The excerpt denotes by \(L(\mathcal A)\) the set of intersections of these hyperplanes, excluding \(\mathbb A^n_K\). It does not specify how empty intersections are treated. Its extracted formula for the flag Hilbert--Poincar\'e series contains notation \(\pi_F(Y)\) and factors involving \(T_X/(1-T_X)\), but the summation range, the meaning of \(F\), and the definition of \(\pi_F(Y)\) are not supplied reliably enough to state the formula precisely.

\emph{Requested construction.} Construct a cohomology theory for subalgebra zeta functions that unifies their arithmetic zeta-function theory with the combinatorial theory of flag Hilbert--Poincar\'e series.

\emph{Unresolved scope.} This is a research-program formulation rather than a fully specified existence statement. The record does not define the class of algebras or subalgebras to be counted, the arithmetic base, the relevant local or global zeta functions, the objects and coefficients of the proposed cohomology theory, or the identities and compatibility properties constituting unification. Moreover, the original question refers to ``these two theories'' without supplying its preceding context; the identification of those theories above follows the supplied cleaned statement. No precise cohomological axioms or recovery formula can therefore be asserted from this record.
\subsection{Short English statement}
Seek a cohomology theory connecting the arithmetic of subalgebra zeta functions with the combinatorics of flag Hilbert--Poincar\'e series. The supplied record leaves the scope and required meaning of this connection unspecified, so this remains a statement-only research question requiring mathematical review.
\subsection{Sources}
\begin{itemize}
\item Supplied local UnsolvedMath registry record, permanent ID \(30006429\), title \emph{Cohomology Theory for Subalgebra Zeta Functions}, source problem code \texttt{catalogue entry}. The record lists proposed year 2025 and supplies no individual proposer.
\item Supplied report citation: \emph{Cohomology Theories for Automorphic Forms and Enumerative Algebra} (2026), Oberwolfach Reports, pp.\ 2065--2124, DOI \texttt{10.4171/owr/2025/38}. \url{https://doi.org/10.4171/owr/2025/38}. The cited year 2026 and the 2025 designation elsewhere in the record are retained without reconciliation. The background describes Oberwolfach Reports as documenting workshops and their open-problem discussions.
\item Supplied literature source, labelled ``Lee--Lee cotype zeta functions'': \url{https://arxiv.org/abs/2504.18908}.
\item Supplied literature source, labelled ``Oberwolfach Report 38'': \url{https://ems.press/content/serial-article-files/52251}.
\item Literature assessment dated exactly \texttt{2026-08-22}: the supplied curation labels the problem open and reports that recent cotype zeta functions generalize and unify several counting zeta functions, while no bespoke cohomology theory unifying the two theories in the extracted question has been established. This is supplied dated curation derived from the OWR record, not an independently verified literature assessment. Its time-sensitive status and cited claims require independent verification.
\end{itemize}
\end{document}
